\documentclass[11pt]{article}
\usepackage[margin=1in]{geometry}
\usepackage{amsmath,amssymb,amsthm}
\usepackage{booktabs}
\usepackage{xcolor}
\usepackage{hyperref}
\hypersetup{colorlinks=true,linkcolor=blue!50!black,urlcolor=blue!50!black,citecolor=blue!50!black}
\usepackage{fancyhdr}
\usepackage[protrusion=true,expansion=false]{microtype}

\pagestyle{fancy}\fancyhf{}
\fancyhead[L]{\small The Relational Representation Principle}
\fancyhead[R]{\small When a Number Lies}
\fancyfoot[C]{\thepage}

\newtheorem{theorem}{Theorem}
\newtheorem{proposition}{Proposition}
\newtheorem{observation}{Observation}
\newtheorem{principle}{Principle}
\newtheorem{definition}{Definition}

\title{\textbf{When a Number Lies:\\
The Relational Representation Principle and Its Limits}\\[4pt]
\large A study of when structural representation beats scalar abstraction}
\author{B.\ Wallace \quad (TensorRent)\\
\small with computational verification by Claude (Anthropic)}
\date{\small Companion to Papers~a--g \ ---\ SIP License v1.1 \ ---\ Seed \texttt{20260423}}

\begin{document}
\maketitle

\begin{abstract}
We isolate and test a single principle: many quantities that computation abstracts to one large scalar are
in fact lossy projections of an exact underlying \emph{relation}, and computing in the relation is more
accurate, sometimes more efficient, and occasionally the only faithful option. We state the principle,
give a representation that realizes it for multiplicative structure (numbers as boxes of prime sides), and
subject it to an adversarial gauntlet across roughly a dozen problems in number theory, combinatorics,
linear algebra, statistical mechanics, and social choice. The principle resolves into a sharp, predictable
law: relational representation wins exactly when the scalar is a lossy-to-impossible projection of an exact
relation that is a \emph{product} (overflow), a \emph{ratio} (cancellation), or a \emph{non-transitive
graph} (no scalar exists); it ties or loses when the scalar is the datum itself, when the relation is
additive, or when an existing method already avoids forming the scalar. We characterize the limits
explicitly. The fundamental wall is addition: a sum destroys multiplicative structure and admits no shape,
proven by direct example. The efficiency claim is sharply bounded: against a compiled big-integer the
relational form does \emph{not} win on speed even at 39{,}000-bit products --- its advantage is overflow
avoidance, exactness, and impossibility, not arithmetic throughput. The strongest instance is non-transitive:
when preferences cycle (Condorcet), \emph{every} scalar ranking provably contradicts the majority relation,
and the cyclicity persists at scale (about $48\%$ of $1001$-voter, five-candidate electorates), so the
impossibility is structural rather than a small-sample artifact. We provide a verified reference engine, a
$10^{80}$-scale verified-optimal structure search as the headline efficiency instance, and a complete record
of the negative results --- including four architecture-integration hypotheses falsified against live source
and two performance hypotheses falsified against measurement --- because the boundary is the contribution.
\end{abstract}

\section{The principle}

\begin{principle}[Relational over scalar]
Let a computation produce a scalar $s$ that is a function of an underlying structured object $R$ (a product,
a ratio, a graph). If $s$ is a lossy or non-invertible projection of $R$, then computing the desired answer
directly on $R$ is more accurate, and often more tractable, than forming $s$ and operating on it.
\end{principle}

The claim is not that scalars are bad. It is that some scalars are \emph{shadows}: the magnitude lives in the
projection, not the object (a $10^{80}$ integer has a small description as an exponent vector; a partition
function with $300$ digits is a short product over modes). When the shadow is what breaks --- overflow,
underflow, catastrophic cancellation, or the outright nonexistence of a faithful scalar --- the object was
the thing to keep. This study asks where that is true, where it is false, and what law separates the two.

\section{The representation and its verification}

\begin{definition}[Shape]
A \emph{shape} encodes a positive integer's multiplicative structure as a finite map $\{p_i:e_i\}$ from
prime sides to positive exponents, with volume $\prod p_i^{e_i}$. Multiplicative questions become geometric:
multiply $=$ merge, divisibility $=$ containment, $\gcd/\mathrm{lcm}=$ common/enclosing box, primality $=$
atomic box. The reference engine is $91$ lines; all twelve operations are verified exact against integer
ground truth (seed \texttt{20260423}). The single boundary is addition, handled by an explicit, flagged exit
to value-space.
\end{definition}

This is the multiplicative realization of the principle. Non-transitive relations require a second
realization --- a directed dominance graph --- developed in \S5.

\section{The gauntlet: where the principle holds}

Each candidate was raced against the \emph{smart} standard method and anchored against ground truth. Results
are graded HIT (relation wins), SPLIT (the multiplicative part wins, the additive part is the wall),
DOWNGRADED (an existing method already wins), or FALSE FIT (no hidden relation).

\begin{center}
\begin{tabular}{lll}
\toprule
domain / problem & grade & mechanism \\
\midrule
non-transitive preference cycles & \textbf{HIT (necessity)} & no scalar ranking can exist \\
highly composite / superabundant search & \textbf{HIT} & exponent-space search; verified-optimal to $10^{80}$ \\
concentrated modular products & \textbf{HIT} & $K/D$ collapse; logs inapplicable to mod \\
exact multiplicities $W=N!/\prod n_i!$ & \textbf{HIT} & logs destroy the integer; relation keeps it \\
lattice index / SNF divisor product & \textbf{HIT} & overflows as scalar; exact as relation \\
partition-function / free-energy ratios & \textbf{HIT} & $Z$ overflows; ratio exact in factor-relation \\
eigenvalue-product determinant & SPLIT & $\det=$ product (hit); trace, char-poly $=$ sums (wall) \\
multiplicative-function families & \textbf{HIT} & whole family free from one held shape \\
content fingerprint & SPLIT & multiplicative: hit; general data: needs a hash \\
probability products & DOWNGRADED & log-space floats already avoid underflow \\
exact rational products & DOWNGRADED & compiled \texttt{Fraction} gcd beats dict merge \\
measured continuous data (temperatures, fluxes) & FALSE FIT & no hidden relation; the scalar is the datum \\
\bottomrule
\end{tabular}
\end{center}

\begin{observation}[The headline efficiency instance, verified-optimal]
Searching directly in exponent space for the integer $n\le N$ maximizing the divisor count $d(n)$ never
forms $n$. An exact branch-and-bound with a \emph{provably admissible} upper bound (budget all remaining
log-room on the cheapest prime, each unit buying at most a factor of two) computes the true maximum out to
$N=10^{80}$ in tens of seconds, sixty-two orders of magnitude past the $\sim 10^{18}$ wall of integer
enumeration. It is anchored to brute force at $N\le 60,840,720720,2162160$ ($d=12,32,240,320$), all exact.
Two earlier versions were \emph{falsified} en route and are retained in the record: a greedy heuristic
(suboptimal, failed the anchor) and a tighter bound that was not admissible (pruned true optima, returned
wrong answers). The verified-optimal version is the one of record.
\end{observation}

\section{The limits (the core contribution)}

A principle without a boundary is a slogan. We pushed to failure and characterized the edge.

\paragraph{The wall: addition.} A sum destroys multiplicative structure and has no shape: $360+84=444=
2^2\cdot 3\cdot 37$ shares no factor structure with either summand. No relational representation crosses
this; it is the fundamental limit, and it reappears in every SPLIT result as the additive component
(trace, characteristic-polynomial coefficients, determinant-as-a-sum).

\paragraph{The efficiency bound.} The principle is \emph{not} an arithmetic accelerator. Forming a product of
$k$ structured factors as a compiled big integer beat the relational representation at every size tested,
through $k=10^4$ ($\sim 39{,}000$ bits). The relational advantage is categorical, not throughput: it avoids
overflow/underflow, preserves exactness, and represents the impossible --- it does not multiply faster.

\paragraph{The requirement.} The principle needs an \emph{exact} underlying relation. Measured continuous
data (a temperature of $2.725$) has no factor structure and no relation to keep; the scalar is the datum and
flattening is correct. Absent an exact relation, the principle does not apply.

\paragraph{The worth-of-exactness bound.} Relational exactness costs time. It is worth paying only when the
question is itself exact or discrete --- modular arithmetic (no approximation exists), exact ratios near
unity (where log-space error would swamp the signal), impossibility cases. When floating tolerance suffices
(most of physics and machine learning), existing methods win and the principle is a refinement at best.

\begin{theorem}[Boundary law]
Relational representation strictly improves on scalar abstraction if and only if (i) an exact underlying
relation exists, (ii) it is a product, a ratio, or a non-transitive graph rather than a sum, and (iii) the
scalar projection breaks (overflow, underflow, cancellation, or nonexistence) or the question is exact and
the scalar destroys needed structure. Otherwise scalar abstraction is correct or already optimally served.
\end{theorem}

This is stated as the organizing law of the study; (i)--(iii) are each witnessed by HITs and each violation
by a DOWNGRADED or FALSE-FIT or SEAM result above. It is an empirical-structural characterization, not a
formal theorem in the proof-theoretic sense.

\section{The strongest instance: when no number can exist}

The deepest form of the principle is not lossiness but impossibility.

\begin{theorem}[Scalar abstraction is impossible under cycles]
If the pairwise majority relation on candidates contains a cycle $A\!\succ\!B\!\succ\!C\!\succ\!A$, then every
total order (every scalar ranking) contradicts at least one majority verdict. The relation carries
information no scalar can hold.
\end{theorem}

\begin{proof}[Verification]
By exhaustive check over all $n!$ rankings on the Condorcet paradox instance, the minimum number of
contradicted majority edges is one, not zero; containment/order is transitive while the majority relation is
not, so no order reproduces it. This is the social-choice instance of the same transitivity obstruction that
forbids encoding rock-paper-scissors as box containment (a partial order cannot represent a cycle).
\end{proof}

\begin{observation}[Structural, not small-sample]
Cyclicity persists and grows with the number of options and does not wash out with electorate size:
measured rates are about $7\%,43\%,78\%,99\%$ of random electorates for $3,5,7,10$ candidates, and about
$48\%$ at $1001$ voters with five candidates. The impossibility is a structural feature of preference
aggregation, not an artifact of few voters.
\end{observation}

This places the principle's strongest claim --- ``some objects cannot be a number, and forcing them loses
truth'' --- on a foundational result (Arrow/Condorcet) rather than on any benchmark.

\section{Negative results retained}

Per the program's discipline, falsifications are first-class outputs. Beyond the two search heuristics
falsified in \S3, this principle's integration into the AISO/OmniForge architecture was tested against live
source and falsified four times: the multiplicative-fit of the BRA charge (the real charge composes
additively by table-sum, confirmed from the kernel whitepaper); a dual-engine geometric router (the live
constraints are ordinal thresholds and the live relation is transitive --- no cycles --- confirmed against
the routing/CDCL source); and trust/Merkle layers (additive and hash-destroyed respectively). The geometry
engines are real, verified, standalone capabilities; they are not that architecture's substrate. The clean
structural resemblance was, in each case, the warning sign, and the source was the arbiter.

\section{Scope and honesty}

The mathematics invoked is classical (fundamental theorem of arithmetic, Legendre's formula, Robin's
criterion connecting superabundant numbers to the Riemann Hypothesis, Arrow/Condorcet, the Hadamard and
Euler products). The contribution is not a theorem but a \emph{characterized representational principle}
with a predictive boundary, witnessed across domains and bounded by explicit limits. The applied target is
accuracy and tractability where scalar abstraction fails; for that target, the principle's established
mathematical components are an asset. The principle does not touch primality or prime-counting (the factoring
wall blocks building a prime's shape), does not prove the Riemann Hypothesis (a categorical gap between finite
multiplicative bookkeeping and the global analytic behavior of an infinite object), and is not an arithmetic
accelerator. What it is: a sharp, testable lens --- \emph{is this scalar hiding a product, a ratio, or a
non-transitive graph that breaks when formed, or is it the truth?} --- with a verified hit set, a verified
miss set, and a fundamental wall at addition.

\vspace{1em}\noindent\rule{\linewidth}{0.4pt}\\
\small Companion to \textit{Colour from Gravity} through \textit{The Geometry Engine} (Papers~a--g). Reference engine, dominance engine, search code, and the domain test
harnesses accompany this study. SIP License v1.1. Seed \texttt{20260423}.

\end{document}
